\slajdid{04-s017}
\begin{frame}[label=04-s017]{Binarni, oktalni i heksadecimalni sistem}
% element: 04-s017:e01 — Binarni, oktalni, heksadecimalni sistemi, osnove i svi skupovi cifara
% element: 04-s017:e02 — Obe dvostrane veze kompatibilnosti
\begin{tabularx}{\linewidth}{@{}llX@{}}\toprule
Sistem & Osnova & Cifre\\\midrule
Binarni & $r=2$ & $c_i\in\{0,1\}$\\[3mm]
Oktalni & $r=8=2^3$ & $c_i\in\{0,1,2,3,4,5,6,7\}$\\[3mm]
Heksadecimalni & $r=16=2^4$ & $c_i\in\{0,1,2,3,4,5,6,7,8,9,A,B,C,D,E,F\}$\\\bottomrule
\end{tabularx}\par\vspace{4mm}
\centering\NaslovBloka{Kompatibilni sistemi}
\begin{tikzpicture}[font=\fontsize{11}{14}\selectfont]
\node (a) at (0,1.2) {binarni};\node (b) at (4.4,1.2) {oktalni};
\node (c) at (0,0) {binarni};\node (d) at (4.4,0) {heksadecimalni};
\draw[<->,DEcrvena,line width=.8pt] ([xshift=2mm]a.east)--([xshift=-2mm]b.west);
\draw[<->,DEcrvena,line width=.8pt] ([xshift=2mm]c.east)--([xshift=-2mm]d.west);
\end{tikzpicture}
\end{frame}
